% ============================================================================
% Rapport PFE - Tablii: Plateforme Multi-tenant de Gestion de Restaurants
% Université TEK-UP -- Diplôme National d'Ingénieur GLSI
% ============================================================================
\documentclass[12pt,a4paper]{report}

% --- Encoding & Language ---
\usepackage[utf8]{inputenc}
\usepackage[T1]{fontenc}
\usepackage[french]{babel}
\usepackage{lmodern}

% --- Page Layout ---
\usepackage[top=2.5cm, bottom=2.5cm, left=2.5cm, right=2.5cm]{geometry}

% --- Colors ---
\usepackage{xcolor}
\definecolor{primary}{RGB}{220,38,38}
\definecolor{secondary}{RGB}{37,99,235}
\definecolor{dark}{RGB}{17,24,39}
\definecolor{lightgray}{RGB}{243,244,246}
\definecolor{accent}{RGB}{16,185,129}

% --- Graphics & Diagrams ---
\usepackage{graphicx}
\usepackage{tikz}
\usetikzlibrary{
    shapes.geometric,
    shapes.multipart,
    arrows.meta,
    positioning,
    calc,
    fit,
    backgrounds,
    decorations.pathreplacing,
    matrix,
    chains,
    shadows,
    patterns,
    trees,
    automata
}

% --- Tables ---
\usepackage{tabularx}
\usepackage{booktabs}
\usepackage{multirow}
\usepackage{longtable}
\usepackage{colortbl}
\usepackage{array}

% --- Code Listings ---
\usepackage{listings}
\lstset{
    basicstyle=\ttfamily\small,
    keywordstyle=\color{secondary}\bfseries,
    commentstyle=\color{gray}\itshape,
    stringstyle=\color{accent},
    numbers=left,
    numberstyle=\tiny\color{gray},
    frame=single,
    rulecolor=\color{lightgray},
    breaklines=true,
    breakatwhitespace=true,
    tabsize=4,
    captionpos=b,
    language=Python
}

% --- Hyperlinks & References ---
\usepackage{hyperref}
\hypersetup{
    colorlinks=true,
    linkcolor=primary,
    urlcolor=secondary,
    citecolor=accent,
    pdfborder={0 0 0}
}

% --- Additional Packages ---
\usepackage{amsmath, amssymb}
\usepackage{enumitem}
\usepackage{fancyhdr}
\usepackage{titlesec}
\usepackage{tocloft}
\usepackage{float}
\usepackage{caption}
\usepackage{subcaption}
\usepackage{makecell}
\usepackage{ragged2e}
\usepackage{setspace}
\usepackage{emptypage}
\usepackage{afterpage}
\usepackage{pifont}
\usepackage{fontawesome5}

% --- Header & Footer ---
\pagestyle{fancy}
\fancyhf{}
\fancyhead[L]{\textcolor{primary}{\leftmark}}
\fancyhead[R]{\textcolor{dark}{Tablii -- Rapport PFE}}
\fancyfoot[C]{\textcolor{gray}{\thepage}}
\renewcommand{\headrulewidth}{1pt}
\renewcommand{\headrule}{\color{primary}\hrule height 1pt}

% --- Section Formatting ---
\titleformat{\chapter}[display]
    {\normalfont\huge\bfseries\color{primary}}
    {\chaptertitlename\ \thechapter}{20pt}{\Huge\color{dark}}
\titleformat{\section}
    {\normalfont\Large\bfseries\color{dark}}
    {\thesection}{1em}{}
\titleformat{\subsection}
    {\normalfont\large\bfseries\color{secondary}}
    {\thesubsection}{1em}{}

% --- Custom Commands ---
\newcommand{\checkmarkitem}{\ding{51}}
\newcommand{\crossitem}{\ding{55}}
\newcommand{\cmark}{\textcolor{accent}{\checkmarkitem}}
\newcommand{\xmark}{\textcolor{primary}{\crossitem}}

\newcommand{\techbadge}[2]{
    \tikz[baseline=(char.base)]{
        \node[char, rectangle, rounded corners=2pt, inner sep=3pt,
              fill=#1!10, draw=#1, text=#1, font=\small\ttfamily] {#2};
    }
}

% --- Line Spacing ---
\onehalfspacing

% ============================================================================
\begin{document}

% ============================================================================
% PAGE DE GARDE
% ============================================================================
\begin{titlepage}
    \centering
    
    % University Header
    \begin{tikzpicture}[remember picture, overlay]
        \fill[primary] (current page.north west) rectangle ++(\paperwidth, -3cm);
        \node[white, font=\Large\bfseries] at (current page.north) {\large
            R\'EPUBLIQUE TUNISIENNE
        };
        \node[white, font=\Large\bfseries] at ([yshift=-1cm]current page.north) {\large
            MINIST\`ERE DE L'ENSEIGNEMENT SUP\'ERIEUR ET DE LA RECHERCHE SCIENTIFIQUE
        };
    \end{tikzpicture}
    
    \vspace{3.5cm}
    
    % University Logo Area
    \begin{tikzpicture}
        \node[draw=primary, line width=2pt, rounded corners, minimum width=6cm, minimum height=2cm, fill=primary!5] (univ) {
            \parbox{5.5cm}{\centering
                \textbf{\Large Universit\'e TEK-UP}\\[0.3cm]
                \textcolor{primary}{\textbf{Ecole d'Ing\'enieurs}}\\
                \textcolor{gray}{\small Sp\'ecialit\'e : GLSI}
            }
        };
    \end{tikzpicture}
    
    \vspace{1.5cm}
    
    % Project Title
    \begin{tikzpicture}
        \fill[primary!10, rounded corners] (-5, -2) rectangle (5, 2);
        \draw[primary, line width=2pt, rounded corners] (-5, -2) rectangle (5, 2);
        \node at (0, 1.2) {\textcolor{gray}{\Large Projet de Fin d'\'Etudes}};
        \node at (0, 0.2) {\textcolor{primary}{\Huge\bfseries Tablii}};
        \node at (0, -0.8) {\textcolor{dark}{\Large\bfseries Plateforme Multi-tenant de}};
        \node at (0, -1.4) {\textcolor{dark}{\Large\bfseries Gestion et Commande de Restaurants}};
    \end{tikzpicture}
    
    \vspace{1.5cm}
    
    % Student & Supervisor Info
    \begin{tabular}{p{6cm} p{6cm}}
        \centering\textbf{\large Pr\'epar\'e par :} & \centering\textbf{\large Encadr\'e par :} \\[0.5cm]
        \centering\textcolor{secondary}{\Large\bfseries Ahmed Trabelsi} & \centering\textcolor{secondary}{\large [Nom Encadrant]} \\
        \centering\textcolor{secondary}{\Large\bfseries Insaffe} & \centering\textcolor{gray}{Enseignant \`a TEK-UP} \\
        \centering\textcolor{secondary}{\Large\bfseries Muslim} & \\
        \centering\textcolor{secondary}{\Large\bfseries Ayoub} & \\[0.5cm]
        \centering\textcolor{gray}{Ing\'enieur en G\'enie Logiciel} & \\
    \end{tabular}
    
    \vspace{2cm}
    
    % Year & Logo
    \begin{tikzpicture}
        \fill[dark, rounded corners] (-3, -0.8) rectangle (3, 0.8);
        \node[white, font=\Large\bfseries] at (0, 0) {Ann\'ee Universitaire 2025--2026};
    \end{tikzpicture}
    
    \vfill
    
    % Bottom bar
    \begin{tikzpicture}[remember picture, overlay]
        \fill[primary] (current page.south west) rectangle ++(\paperwidth, 1.5cm);
        \coordinate (bl) at (current page.south west);
        \node[white, font=\small] at ($(bl) + (\paperwidth/2, 0.75cm)$) {
            \textbf{Tablii} -- Menu QR Code, Commandes en Temps R\'eel, Gestion Multi-Restaurant
        };
    \end{tikzpicture}
\end{titlepage}

% ============================================================================
% DEDICACE
% ============================================================================
\cleardoublepage
\chapter*{D\'edicaces}
\addcontentsline{toc}{chapter}{D\'edicaces}

\begin{center}
    \vspace{2cm}
    \textit{\large
    A nos parents, pour leur soutien inconditionnel et leurs sacrifices.\\[0.5cm]
    A nos enseignants, pour leur guidance pr\'ecieuse tout au long de notre parcours.\\[0.5cm]
    A tous ceux qui croient en la puissance de la technologie\\
    pour transformer l'exp\'erience culinaire en Tunisie.
    }
    \vspace{2cm}
\end{center}

% ============================================================================
% REMERCIEMENTS
% ============================================================================
\cleardoublepage
\chapter*{Remerciements}
\addcontentsline{toc}{chapter}{Remerciements}

Nous tenons \`a exprimer notre profonde gratitude \`a :

\begin{itemize}[leftmargin=2cm, itemsep=0.5cm]
    \item Notre encadrant, pour son accompagnement, ses conseils avis\'es et sa disponibilit\'e tout au long de ce projet.
    \item L'\'equipe p\'edagogique de TEK-UP, pour la qualit\'e de la formation dispens\'ee durant notre cycle d'ing\'enieur.
    \item Tous les professionnels de la restauration tunisienne qui ont contribu\'e \`a l'analyse des besoins et \`a la validation de notre solution.
    \item Nos coll\`egues et amis, pour leur soutien moral et technique durant les phases de d\'eveloppement.
\end{itemize}

% ============================================================================
% RESUME
% ============================================================================
\cleardoublepage
\chapter*{R\'esum\'e}
\addcontentsline{toc}{chapter}{R\'esum\'e}

\begin{center}
    \fbox{\parbox{0.9\textwidth}{
        \textbf{Mots-cl\'es :} Restaurant, Multi-tenant, QR Code, WebSocket, Flask, Python, 
        Temps r\'eel, SaaS, Paiement en ligne, PWA, Analytics, Fid\'elisation, Tunisie
    }}
\end{center}

\vspace{0.5cm}

\textbf{R\'esum\'e :}

Tablii est une plateforme SaaS multi-tenant con\c{c}ue pour digitaliser compl\`etement 
l'exp\'erience de restauration en Tunisie. Elle permet aux clients de scanner un QR code sur leur table, 
consulter le menu multilingue (fran\c{c}ais, arabe, anglais) et passer des commandes directement 
depuis leur smartphone, sans installation d'application.

L'architecture repose sur Flask (Python) avec SQLAlchemy comme ORM, Flask-SocketIO pour les 
communications temps r\'eel, et Tailwind CSS pour l'interface utilisateur. La plateforme supporte 
plusieurs restaurants ind\'ependants (multi-tenant), chacun avec son propre menu, son personnel 
(caissier, cuisine, serveur), ses abonnements et ses statistiques.

Les fonctionnalit\'es cl\'es incluent : les commandes en temps r\'eel via WebSocket, un tableau 
de bord analytique avec heatmaps et tendances, l'int\'egration de paiement Flouci (passerelle 
tunisienne), un syst\`eme de fid\'elisation, un mode Ramadan (Iftar/Suhoor), et un support PWA 
pour une exp\'erience mobile optimale.

\vspace{0.5cm}

\noindent\textbf{Abstract :}

Tablii is a multi-tenant SaaS platform designed to fully digitize the restaurant experience in Tunisia. 
It enables customers to scan a QR code on their table, browse a multilingual menu (French, 
Arabic, English), and place orders directly from their smartphone without installing any app.

The architecture is built on Flask (Python) with SQLAlchemy as ORM, Flask-SocketIO for 
real-time communications, and Tailwind CSS for the user interface. The platform supports 
multiple independent restaurants (multi-tenant), each with its own menu, staff (cashier, 
kitchen, waiter), subscriptions, and analytics.

Key features include: real-time orders via WebSocket, an analytics dashboard with heatmaps 
and trends, Flouci payment gateway integration (Tunisian), a loyalty system, Ramadan mode 
(Iftar/Suhoor), and PWA support for an optimal mobile experience.

% ============================================================================
% TABLE DES MATIERES
% ============================================================================
\cleardoublepage
\tableofcontents
\cleardoublepage
\listoffigures
\addcontentsline{toc}{chapter}{Liste des Figures}
\cleardoublepage
\listoftables
\addcontentsline{toc}{chapter}{Liste des Tables}

% ============================================================================
% LISTE DES ACRONYMES
% ============================================================================
\cleardoublepage
\chapter*{Liste des Acronymes}
\addcontentsline{toc}{chapter}{Liste des Acronymes}

\begin{tabularx}{\textwidth}{l X}
    \toprule
    \textbf{Acronyme} & \textbf{Signification} \\
    \midrule
    SaaS & Software as a Service \\
    PWA & Progressive Web Application \\
    API & Application Programming Interface \\
    ORM & Object-Relational Mapping \\
    WebSocket & Protocole de communication bidirectionnelle temps r\'eel \\
    QR Code & Quick Response Code \\
    CSRF & Cross-Site Request Forgery \\
    ER & Entity-Relationship \\
    MVP & Minimum Viable Product \\
    CI/CD & Continuous Integration / Continuous Deployment \\
    KDS & Kitchen Display System \\
    TND & Dinar Tunisien \\
    RTL & Right-To-Left (pour l'arabe) \\
    \bottomrule
\end{tabularx}

% ============================================================================
% INTRODUCTION GENERALE
% ============================================================================
\cleardoublepage
\chapter*{Introduction G\'en\'erale}
\addcontentsline{toc}{chapter}{Introduction G\'en\'erale}

L'essor des technologies numériques a profondément transformé les modèles économiques traditionnels, et le secteur de la restauration en Tunisie ne fait pas exception à cette règle. Avec la démocratisation des smartphones et l'amélioration continue de la connectivité internet, les attentes des consommateurs en matière de rapidité, de transparence et d'interactivité sont en constante évolution. Pourtant, force est de constater que la gestion quotidienne d'un nombre significatif de restaurants et cafés tunisiens demeure encore ancrée dans des processus artisanaux : menus papier sujets à l'usure et difficiles à mettre à jour, prise de commande manuelle génératrice d'erreurs de transmission entre les serveurs et la cuisine, et absence de visibilité en temps réel sur l'activité commerciale pour les gérants.

Ces limitations opérationnelles constituent un frein majeur à la scalabilité et à l'optimisation de la rentabilité des établissements. La numérisation de ces processus n'est plus une simple option de confort, mais un impératif stratégique pour rester compétitif dans un marché de plus en plus exigeant. C'est dans ce contexte que s'inscrit notre projet de fin d'études, intitulé \textbf{Tablii}.

L'objectif principal de ce projet est de concevoir et de développer une application SaaS (Software as a Service) multi-tenant dédiée à la gestion intelligente des restaurants et cafés. Tablii propose une solution intégrée permettant aux clients de consulter le menu digital et de passer commande directement via un QR Code, tout en offrant aux équipes de service (caissiers, serveurs et personnel de cuisine) un dashboard en temps réel pour le suivi des commandes et des tables. Parallèlement, un panel d'administration robuste offre aux gérants la capacité de gérer leurs établissements, leurs menus et d'analyser les statistiques de vente avec une granularité fine.

Sur le plan technique, le choix d'architecture a été guidé par la nécessité d'une performance optimale et d'une accessibilité universelle. Le frontend a été développé en \textbf{Vanilla JavaScript} couplé à \textbf{Tailwind CSS}, garantissant une légèreté exemplaire et une fluidité d'exécution même sur les appareils mobiles d'entrée de gamme souvent utilisés dans les environnements de restauration, tout en évitant la surcharge liée aux frameworks lourds. Le backend, basé sur \textbf{Flask (Python)}, assure une gestion rigoureuse du multi-tenancy, permettant à plusieurs établissements de coexister sur la même infrastructure avec une isolation stricte des données.

Le présent rapport s'articule autour de quatre chapitres principaux :
\begin{itemize}
    \item Le \textbf{Chapitre 1} dresse le cadre général du projet, présente l'étude de l'existant et justifie la solution proposée.
    \item Le \textbf{Chapitre 2} détaille l'analyse et la spécification des besoins fonctionnels et non-fonctionnels, illustrés par les diagrammes UML pertinents.
    \item Le \textbf{Chapitre 3} expose la conception architecturale de l'application, incluant le modèle de données multi-tenant et les diagrammes de conception.
    \item Le \textbf{Chapitre 4} est consacré à la réalisation et à l'implémentation, mettant en avant les choix techniques, les interfaces développées et les résultats obtenus.
\end{itemize}

Enfin, une conclusion générale viendra synthétiser les apports de ce travail et ouvrir sur les perspectives d'évolution de Tablii.

% ============================================================================
% CHAPITRE 1 : CONTEXTE GENERAL DU PROJET
% ============================================================================
\cleardoublepage
\chapter{Contexte Général du Projet}

\section{Introduction}

Ce chapitre présente le cadre général du projet Tablii, le contexte de la restauration en Tunisie, une étude critique des solutions existantes sur le marché, et la justification de notre solution proposée.

\section{Présentation du Cadre du Projet}

\subsection{Contexte du Secteur de la Restauration en Tunisie}

Le secteur de la restauration en Tunisie connaît une croissance soutenue, avec une densité élevée de restaurants, cafés et salons de thé dans les zones urbaines. Cependant, la majorité de ces établissements fonctionnent encore avec des méthodes de gestion traditionnelles :

\begin{itemize}[leftmargin=1.5cm, itemsep=0.3cm]
    \item \textbf{Menus papier} : Coûteux à imprimer, difficiles à mettre à jour en temps réel (changement de prix, rupture de stock), et non adaptés aux touristes étrangers.
    \item \textbf{Prise de commande manuelle} : Le serveur note la commande sur un carnet, la transmet oralement ou physiquement à la cuisine, ce qui génère des erreurs de communication et des pertes de temps.
    \item \textbf{Suivi de commande inexistant} : Le client n'a aucune visibilité sur l'état de sa commande, ce qui crée de l'anxiété et des demandes répétées auprès du serveur.
    \item \textbf{Gestion financière artisanale} : Le calcul des ventes, des statistiques et des tendances se fait manuellement, rendant l'analyse commerciale fastidieuse et imprécise.
\end{itemize}

\subsection{Objectifs du Projet}

Le projet Tablii vise à répondre à ces problématiques en offrant une plateforme complète qui :

\begin{enumerate}[leftmargin=1.5cm, itemsep=0.3cm]
    \item \textbf{Digitalise le menu} : Remplace le menu papier par un menu digital accessible via QR Code, multilingue (FR/AR/EN) et mis à jour en temps réel.
    \item \textbf{Automatise la prise de commande} : Le client commande directement depuis son smartphone, éliminant les erreurs de transmission.
    \item \textbf{Connecte tous les acteurs} : Synchronise en temps réel le client, le caissier, la cuisine et le serveur via WebSocket.
    \item \textbf{Fournit des analytics} : Offre au gérant un tableau de bord avec des statistiques de vente, des heatmaps d'heures de pointe et des articles populaires.
    \item \textbf{Supporte le multi-tenant} : Permet à plusieurs restaurants de coexister sur la même plateforme avec isolation complète des données.
    \item \textbf{Intègre le paiement local} : Connecte la passerelle de paiement tunisienne Flouci pour le paiement en ligne.
\end{enumerate}

\section{Étude de l'Existant}

\subsection{Solutions Concurrentes sur le Marché}

Nous avons analysé trois solutions majeures présentes sur le marché de la restauration digitale :

\subsubsection{GloriaFood}

GloriaFood est une plateforme de commande en ligne qui propose un système de menu digital et de prise de commande en ligne.

\begin{table}[H]
    \centering
    \begin{tabularx}{\textwidth}{l X}
        \toprule
        \rowcolor{primary!10}
        \textbf{Points Forts} & \textbf{Points Faibles} \\
        \midrule
        Interface intuitive & Pas de temps réel (pas de WebSocket) \\
        Système de commande en ligne mature & Pas de support multilingue AR/FR/EN natif \\
        Application mobile disponible & Pas adapté aux spécificités tunisiennes \\
        \bottomrule
    \end{tabularx}
    \caption{Analyse de GloriaFood}
    \label{tab:gloriafood}
\end{table}

\subsubsection{Toast POS}

Toast POS est un système de point de vente complet pour les restaurants, incluant la gestion des commandes, des paiements et des inventaires.

\begin{table}[H]
    \centering
    \begin{tabularx}{\textwidth}{l X}
        \toprule
        \rowcolor{secondary!10}
        \textbf{Points Forts} & \textbf{Points Faibles} \\
        \midrule
        Temps réel intégré & Solution coûteuse (abonnement mensuel élevé) \\
        Gestion complète du restaurant & Pas de multi-tenant SaaS (solution on-premise) \\
        Matériel dédié disponible & Pas de support pour le marché tunisien \\
        \bottomrule
    \end{tabularx}
    \caption{Analyse de Toast POS}
    \label{tab:toast}
\end{table}

\subsubsection{Square for Restaurants}

Square propose une solution de gestion de restaurant avec point de vente, menu digital et analytics.

\begin{table}[H]
    \centering
    \begin{tabularx}{\textwidth}{l X}
        \toprule
        \rowcolor{accent!10}
        \textbf{Points Forts} & \textbf{Points Faibles} \\
        \midrule
        Écosystème complet (paiement + gestion) & Pas de multi-tenant SaaS \\
        Analytics avancés & Pas de support RTL pour l'arabe \\
        Intégration livraison & Pas de mode Ramadan ou adaptation locale \\
        \bottomrule
    \end{tabularx}
    \caption{Analyse de Square for Restaurants}
    \label{tab:square}
\end{table}

\subsection{Synthèse Comparative}

\begin{table}[H]
    \centering
    \begin{tabularx}{\textwidth}{l c c c c}
        \toprule
        \rowcolor{primary!10}
        \textbf{Critère} & \textbf{GloriaFood} & \textbf{Toast POS} & \textbf{Square} & \textbf{Tablii} \\
        \midrule
        Menu QR Code & \cmark & \cmark & \cmark & \cmark \\
        Temps réel (WebSocket) & \xmark & \cmark & \cmark & \cmark \\
        Multi-tenant SaaS & \cmark & \xmark & \xmark & \cmark \\
        Multilingue AR/FR/EN & \xmark & \xmark & \xmark & \cmark \\
        Paiement local (Tunisie) & \xmark & \xmark & \xmark & \cmark \\
        Mode Ramadan & \xmark & \xmark & \xmark & \cmark \\
        PWA & \xmark & \xmark & \xmark & \cmark \\
        Prix abordable & \cmark & \xmark & \xmark & \cmark \\
        \bottomrule
    \end{tabularx}
    \caption{Comparaison synthétique des solutions}
    \label{tab:comparison}
\end{table}

\subsection{Critique de l'Existant et Opportunité}

L'analyse des solutions existantes révèle plusieurs lacunes majeures :

\begin{itemize}[leftmargin=1.5cm, itemsep=0.3cm]
    \item \textbf{Absence d'adaptation locale} : Aucune solution ne supporte le multilingue arabe-français-anglais avec le RTL, ni le mode Ramadan (Iftar/Suhoor).
    \item \textbf{Coût élevé} : Les solutions complètes comme Toast et Square sont financièrement inaccessibles pour les petits et moyens restaurants tunisiens.
    \item \textbf{Pas de paiement local} : Aucune intégration avec les passerelles de paiement tunisiennes (Flouci, Konnect, etc.).
    \item \textbf{Architecture non-SaaS} : La plupart des solutions ne sont pas multi-tenant, ce qui rend la gestion de plusieurs restaurants complexe.
\end{itemize}

C'est dans ce contexte que \textbf{Tablii} se positionne comme une solution innovante, abordable et adaptée aux spécificités du marché tunisien.

\section{Solution Proposée : Tablii}

Tablii est une plateforme SaaS multi-tenant qui digitalise l'ensemble du parcours de restauration :

\begin{figure}[H]
    \centering
    \begin{tikzpicture}[
        box/.style={rectangle, rounded corners, draw=#1, fill=#1!5, minimum width=3.5cm, minimum height=1cm, font=\small\bfseries, text=#1},
        arrow/.style={-{Stealth}, thick, color=gray}
    ]
        \node[box=primary] (client) at (0, 4) {Client (QR Code)};
        \node[box=secondary] (cashier) at (-5, 1) {Caissier};
        \node[box=accent] (kitchen) at (0, 1) {Cuisine};
        \node[box=orange] (waiter) at (5, 1) {Serveur};
        \node[box=purple] (owner) at (-3, -2) {Gérant};
        \node[box=teal] (admin) at (3, -2) {Super Admin};
        
        \draw[arrow] (client) -- (cashier);
        \draw[arrow] (client) -- (kitchen);
        \draw[arrow] (client) -- (waiter);
        \draw[arrow] (cashier) -- (kitchen);
        \draw[arrow] (waiter) -- (kitchen);
        \draw[arrow] (owner) -- (cashier);
        \draw[arrow] (owner) -- (kitchen);
        \draw[arrow] (owner) -- (waiter);
        \draw[arrow] (admin) -- (owner);
    \end{tikzpicture}
    \caption{Architecture fonctionnelle de Tablii}
    \label{fig:functional-arch}
\end{figure}

\section{Méthodologie de Gestion de Projet}

Pour mener à bien ce projet, nous avons adopté la méthodologie \textbf{Scrum}, une approche agile qui permet une itération rapide et une adaptation continue aux besoins.

\subsection{Principes de Scrum}

Scrum se base sur des sprints de durée fixe (2 à 3 semaines), avec des rôles bien définis :

\begin{itemize}[leftmargin=1.5cm, itemsep=0.3cm]
    \item \textbf{Product Owner} : Définit les priorités du Product Backlog.
    \item \textbf{Scrum Master} : Facilite le processus et remove les obstacles.
    \item \textbf{Development Team} : Réalise les tâches techniques du Sprint Backlog.
\end{itemize}

\subsection{Organisation de Notre Projet}

Notre projet a été structuré en \textbf{12 phases} (sprints) couvrant l'ensemble du cycle de développement :

\begin{table}[H]
    \centering
    \begin{tabularx}{\textwidth}{c l X}
        \toprule
        \rowcolor{primary!10}
        \textbf{Phase} & \textbf{Sprint} & \textbf{Objectif} \\
        \midrule
        1 & Fondation & Configuration du projet, Flask, Tailwind CSS \\
        2 & Modèles & Base de données multi-tenant (SQLAlchemy) \\
        3 & Utilitaires & Helpers, upload d'images, QR Code \\
        4 & Authentification & Login/Register pour owners et staff \\
        5 & Interface Client & Menu QR Code, panier, commande \\
        6 & Dashboard & Analytics, gestion du menu et des tables \\
        7 & Interfaces Staff & Caissier, Cuisine, Serveur \\
        8 & WebSocket & Temps réel avec Flask-SocketIO \\
        9 & Paiement & Intégration Flouci, webhooks \\
        10 & PWA & Service Worker, manifest, offline \\
        11 & API \& Tests & API REST, tests pytest, coverage \\
        12 & Déploiement & Render.com, PostgreSQL, CI/CD \\
        \bottomrule
    \end{tabularx}
    \caption{Planification des sprints du projet Tablii}
    \label{tab:sprints}
\end{table}

% ============================================================================
% CHAPITRE 2 : ANALYSE ET SPECIFICATION DES BESOINS
% ============================================================================
\cleardoublepage
\chapter{Analyse et Spécification des Besoins}

\section{Introduction}

Ce chapitre détaille l'analyse des besoins fonctionnels et non-fonctionnels du système Tablii, en identifiant les acteurs, les cas d'utilisation et les exigences techniques.

\section{Identification des Acteurs}

Le système Tablii implique six acteurs principaux, chacun avec des besoins et des permissions spécifiques :

\begin{figure}[H]
    \centering
    \begin{tikzpicture}[
        actor/.style={circle, draw=#1, fill=#1!10, minimum size=1.5cm, font=\small\bfseries, text=#1},
        desc/.style={rectangle, rounded corners, draw=gray!50, fill=gray!5, text width=4cm, inner sep=8pt, font=\small}
    ]
        % Actors
        \node[actor=primary] (client) at (0, 5) {Client};
        \node[actor=secondary] (cashier) at (-5, 2) {Caissier};
        \node[actor=accent] (kitchen) at (0, 2) {Cuisine};
        \node[actor=orange] (waiter) at (5, 2) {Serveur};
        \node[actor=purple] (owner) at (-3.5, -1.5) {Gérant};
        \node[actor=teal] (admin) at (3.5, -1.5) {Super Admin};
        
        % Descriptions
        \node[desc, right=0.3cm] at (1.5, 5) {
            Consulte le menu via QR code, passe des commandes, paie, laisse un avis
        };
        \node[desc, left=0.3cm] at (-6.8, 2) {
            Reçoit et valide les commandes, gère les paiements en caisse
        };
        \node[desc, right=0.3cm] at (1.8, 2) {
            Voit les commandes en temps réel, met à jour le statut de préparation
        };
        \node[desc, right=0.3cm] at (6.8, 2) {
            Reçoit les appels, sert les plats, gère les tables assignées
        };
        \node[desc, below=0.3cm] at (-3.5, -2.5) {
            Gère le menu, les tables, le personnel, consulte les statistiques
        };
        \node[desc, below=0.3cm] at (3.5, -2.5) {
            Supervise tous les restaurants, gère les abonnements et la plateforme
        };
        
        % Connections
        \draw[->, thick, gray!50] (client) -- (cashier);
        \draw[->, thick, gray!50] (client) -- (kitchen);
        \draw[->, thick, gray!50] (client) -- (waiter);
        \draw[->, thick, gray!50] (cashier) -- (kitchen);
        \draw[->, thick, gray!50] (waiter) -- (kitchen);
        \draw[->, thick, gray!50] (owner) -- (cashier);
        \draw[->, thick, gray!50] (owner) -- (kitchen);
        \draw[->, thick, gray!50] (owner) -- (waiter);
        \draw[->, thick, gray!50] (admin) -- (owner);
    \end{tikzpicture}
    \caption{Diagramme des acteurs du système Tablii}
    \label{fig:actors}
\end{figure}

\section{Spécification des Besoins Fonctionnels}

\subsection{Besoins Fonctionnels du Client}

\begin{table}[H]
    \centering
    \begin{tabularx}{\textwidth}{c l X}
        \toprule
        \rowcolor{primary!10}
        \textbf{ID} & \textbf{Fonctionnalité} & \textbf{Description} \\
        \midrule
        BF-C01 & Scan QR Code & Accéder au menu en scannant le QR code de la table \\
        BF-C02 & Menu multilingue & Consulter le menu en français, arabe ou anglais \\
        BF-C03 & Passer commande & Sélectionner des articles et valider la commande \\
        BF-C04 & Personnalisation & Choisir des options (taille, suppléments, etc.) \\
        BF-C05 & Appeler serveur & Demander l'assistance d'un serveur \\
        BF-C06 & Paiement en ligne & Payer via Flouci (carte bancaire) \\
        BF-C07 & Laisser un avis & Noter et commenter après la commande \\
        BF-C08 & Programme fidélité & Accumuler et utiliser des points de fidélité \\
        \bottomrule
    \end{tabularx}
    \caption{Besoins fonctionnels -- Client}
    \label{tab:bf-client}
\end{table}

\subsection{Besoins Fonctionnels du Personnel}

\begin{table}[H]
    \centering
    \begin{tabularx}{\textwidth}{c l l X}
        \toprule
        \rowcolor{secondary!10}
        \textbf{ID} & \textbf{Rôle} & \textbf{Fonctionnalité} & \textbf{Description} \\
        \midrule
        BF-S01 & \multirow{3}{*}{Caissier} & Voir commandes & Voir les commandes entrantes en temps réel \\
        BF-S02 & & Accepter/Refuser & Valider ou rejeter une commande \\
        BF-S03 & & Encaisser & Gérer les paiements en caisse \\
        \midrule
        BF-S04 & \multirow{3}{*}{Cuisine} & Voir commandes & Voir les commandes acceptées \\
        BF-S05 & & Changer statut & Passer de "en préparation" à "prêt" \\
        BF-S06 & & Notifications & Recevoir les nouvelles commandes via WebSocket \\
        \midrule
        BF-S07 & \multirow{3}{*}{Serveur} & Voir appels & Voir les appels des clients \\
        BF-S08 & & Servir plats & Confirmer la livraison des plats \\
        BF-S09 & & Gérer tables & Voir les tables assignées et leur statut \\
        \bottomrule
    \end{tabularx}
    \caption{Besoins fonctionnels -- Personnel de restaurant}
    \label{tab:bf-staff}
\end{table}

\subsection{Besoins Fonctionnels du Gérant et Super Admin}

\begin{table}[H]
    \centering
    \begin{tabularx}{\textwidth}{c l l X}
        \toprule
        \rowcolor{accent!10}
        \textbf{ID} & \textbf{Rôle} & \textbf{Fonctionnalité} & \textbf{Description} \\
        \midrule
        BF-A01 & \multirow{5}{*}{Gérant} & Dashboard & Vue d'ensemble des statistiques \\
        BF-A02 & & Gérer menu & Ajouter/modifier/supprimer des articles \\
        BF-A03 & & Gérer tables & Configurer les tables et les QR codes \\
        BF-A04 & & Gérer staff & Créer et gérer les comptes du personnel \\
        BF-A05 & & Analytics & Voir les revenus, heatmaps, tendances \\
        \midrule
        BF-A06 & \multirow{3}{*}{Super Admin} & Tous restaurants & Superviser tous les restaurants inscrits \\
        BF-A07 & & Abonnements & Gérer les plans et les paiements \\
        BF-A08 & & Statistiques globales & Vue consolidée de la plateforme \\
        \bottomrule
    \end{tabularx}
    \caption{Besoins fonctionnels -- Gérant \& Super Admin}
    \label{tab:bf-admin}
\end{table}

\section{Spécification des Besoins Non Fonctionnels}

\begin{table}[H]
    \centering
    \begin{tabularx}{\textwidth}{c l X}
        \toprule
        \rowcolor{dark!10}
        \textbf{ID} & \textbf{Catégorie} & \textbf{Description} \\
        \midrule
        BNF-01 & Performance & Temps de réponse $<$ 200ms pour les notifications WebSocket \\
        BNF-02 & Scalabilité & Support de 50+ restaurants simultanés sur la même instance \\
        BNF-03 & Sécurité & Protection CSRF, hashage bcrypt, sessions sécurisées \\
        BNF-04 & Disponibilité & 99.9\% uptime en production \\
        BNF-05 & Multi-langue & Support FR/AR/EN avec RTL pour l'arabe \\
        BNF-06 & Responsive & Compatible mobile, tablette et desktop \\
        BNF-07 & PWA & Installable, fonctionnement offline partiel \\
        BNF-08 & Multi-tenant & Isolation complète des données entre restaurants \\
        BNF-09 & Légèreté & Frontend en Vanilla JS + Tailwind CSS (pas de framework lourd) \\
        BNF-10 & Compatibilité & Fonctionne sur les appareils d'entrée de gamme \\
        \bottomrule
    \end{tabularx}
    \caption{Besoins non fonctionnels}
    \label{tab:bnf}
\end{table}

\section{Diagramme de Cas d'Utilisation Global}

\begin{figure}[H]
    \centering
    \begin{tikzpicture}[scale=0.85,
        usecase/.style={ellipse, draw=secondary, fill=secondary!5, minimum width=3.5cm, minimum height=1cm, inner sep=5pt, font=\small},
        actor/.style={font=\small\bfseries},
        system/.style={draw=primary, thick, rounded corners, fill=primary!2}
    ]
        % System boundary
        \node[system, minimum width=14cm, minimum height=12cm] (system) at (0, 0) {};
        \node[above=0.2cm] at (system.north) {\textbf{\large Système Tablii}};
        
        % Use cases
        \node[usecase] (uc1) at (-4, 4) {Scanner QR Code};
        \node[usecase] (uc2) at (4, 4) {Consulter le Menu};
        \node[usecase] (uc3) at (-5, 2) {Passer Commande};
        \node[usecase] (uc4) at (5, 2) {Personnaliser un Plat};
        \node[usecase] (uc5) at (-4, 0) {Payer en Ligne};
        \node[usecase] (uc6) at (4, 0) {Appeler Serveur};
        \node[usecase] (uc7) at (-5, -2) {Laisser un Avis};
        \node[usecase] (uc8) at (5, -2) {Gérer le Menu};
        \node[usecase] (uc9) at (-4, -4) {Dashboard Analytics};
        \node[usecase] (uc10) at (4, -4) {Gérer les Commandes};
        
        % Actors
        \node[actor] (client) at (-9, 3) {Client};
        \node[actor] (staff) at (9, 1) {Personnel};
        \node[actor] (owner) at (-9, -3) {Gérant};
        \node[actor] (admin) at (9, -3) {Super Admin};
        
        % Connections - Client
        \draw[->] (client) -- (uc1);
        \draw[->] (client) -- (uc2);
        \draw[->] (client) -- (uc3);
        \draw[->] (client) -- (uc5);
        \draw[->] (client) -- (uc6);
        \draw[->] (client) -- (uc7);
        
        % Connections - Staff
        \draw[->] (staff) -- (uc10);
        
        % Connections - Owner
        \draw[->] (owner) -- (uc8);
        \draw[->] (owner) -- (uc9);
        
        % Connections - Admin
        \draw[->] (admin) -- (uc9);
        
        % Include/Extend
        \draw[dashed, ->] (uc3) -- node[above, font=\small] {\textit{<<include>>}} (uc4);
        \draw[dashed, ->] (uc5) -- node[below, font=\small] {\textit{<<extend>>}} (uc3);
    \end{tikzpicture}
    \caption{Diagramme de cas d'utilisation global}
    \label{fig:usecase}
\end{figure}

\section{Diagramme de Classes Global}

\begin{figure}[H]
    \centering
    \begin{tikzpicture}[scale=0.7,
        entity/.style={rectangle, draw=secondary, fill=secondary!5, minimum width=3cm, minimum height=1.5cm, font=\small\bfseries},
        attr/.style={font=\tiny, text=gray},
        rel/.style={-{Stealth}, thick, gray}
    ]
        % Users
        \node[entity] (users) at (-10, 6) {users};
        \node[attr, below=0.1cm] at (users.south) {id, email, password\_hash, name, role};
        
        % Staff
        \node[entity] (staff) at (-5, 6) {staff\_users};
        \node[attr, below=0.1cm] at (staff.south) {id, restaurant\_id, username, role};
        
        % Restaurant
        \node[entity, fill=primary!5, draw=primary] (restaurant) at (0, 6) {restaurants};
        \node[attr, below=0.1cm] at (restaurant.south) {id, owner\_id, name, slug, city};
        
        % Subscription
        \node[entity] (subscription) at (5, 6) {subscriptions};
        \node[attr, below=0.1cm] at (subscription.south) {id, restaurant\_id, plan, max\_tables};
        
        % Categories
        \node[entity] (categories) at (-7, 2) {categories};
        \node[attr, below=0.1cm] at (categories.south) {id, restaurant\_id, name\_fr, name\_ar, name\_en};
        
        % Menu Items
        \node[entity] (menu_items) at (-2, 2) {menu\_items};
        \node[attr, below=0.1cm] at (menu_items.south) {id, category\_id, price, is\_available};
        
        % Customizations
        \node[entity] (custom) at (3, 2) {customizations};
        \node[attr, below=0.1cm] at (custom.south) {id, menu\_item\_id, selection\_type};
        
        % Tables
        \node[entity] (tables) at (8, 2) {tables\_};
        \node[attr, below=0.1cm] at (tables.south) {id, restaurant\_id, table\_number, qr\_code\_url};
        
        % Table Sessions
        \node[entity] (sessions) at (8, -2) {table\_sessions};
        \node[attr, below=0.1cm] at (sessions.south) {id, table\_id, session\_token, guest\_name};
        
        % Orders
        \node[entity, fill=accent!5, draw=accent] (orders) at (0, -2) {orders};
        \node[attr, below=0.1cm] at (orders.south) {id, session\_id, status, total\_amount};
        
        % Order Items
        \node[entity] (order_items) at (-5, -2) {order\_items};
        \node[attr, below=0.1cm] at (order_items.south) {id, order\_id, menu\_item\_id, quantity};
        
        % Payments
        \node[entity] (payments) at (5, -2) {payment\_transactions};
        \node[attr, below=0.1cm] at (payments.south) {id, order\_id, gateway, status};
        
        % Reviews
        \node[entity] (reviews) at (-10, -2) {reviews};
        \node[attr, below=0.1cm] at (reviews.south) {id, order\_id, rating, comment};
        
        % Customers
        \node[entity] (customers) at (-10, 0) {customers};
        \node[attr, below=0.1cm] at (customers.south) {id, phone, name};
        
        % Loyalty
        \node[entity] (loyalty) at (-10, -4) {loyalty\_points};
        \node[attr, below=0.1cm] at (loyalty.south) {id, customer\_id, restaurant\_id, points};
        
        % Relationships
        \draw[rel] (users) -- node[above, font=\tiny] {1:N} (restaurant);
        \draw[rel] (restaurant) -- node[above, font=\tiny] {1:N} (staff);
        \draw[rel] (restaurant) -- node[above, font=\tiny] {1:1} (subscription);
        \draw[rel] (restaurant) -- node[left, font=\tiny] {1:N} (categories);
        \draw[rel] (restaurant) -- node[right, font=\tiny] {1:N} (tables);
        \draw[rel] (restaurant) -- node[right, font=\tiny] {1:N} (orders);
        \draw[rel] (categories) -- node[above, font=\tiny] {1:N} (menu_items);
        \draw[rel] (menu_items) -- node[above, font=\tiny] {1:N} (custom);
        \draw[rel] (tables) -- node[right, font=\tiny] {1:N} (sessions);
        \draw[rel] (sessions) -- node[above, font=\tiny] {1:N} (orders);
        \draw[rel] (orders) -- node[above, font=\tiny] {1:N} (order_items);
        \draw[rel] (orders) -- node[above, font=\tiny] {1:1} (payments);
        \draw[rel] (orders) -- node[above, font=\tiny] {1:1} (reviews);
        \draw[rel] (customers) -- node[right, font=\tiny] {1:N} (sessions);
        \draw[rel] (customers) -- node[right, font=\tiny] {1:N} (loyalty);
        \draw[rel] (menu_items) -- node[left, font=\tiny] {1:N} (order_items);
    \end{tikzpicture}
    \caption{Diagramme de classes global simplifié}
    \label{fig:class-diagram}
\end{figure}

% ============================================================================
% CHAPITRE 3 : CONCEPTION ARCHITECTURELLE
% ============================================================================
\cleardoublepage
\chapter{Conception Architecturelle}

\section{Introduction}

Ce chapitre expose l'architecture globale de Tablii, le modèle de données multi-tenant, les diagrammes de séquence et les choix de conception qui sous-tendent l'implémentation.

\section{Architecture Globale}

Tablii suit une architecture \textbf{MVC (Modèle-Vue-Contrôleur)} avec une couche de services métier séparée et des WebSockets pour le temps réel.

\begin{figure}[H]
    \centering
    \begin{tikzpicture}[
        layer/.style={rectangle, rounded corners, draw, fill=#1!10, minimum width=11cm, minimum height=1.2cm, font=\small\bfseries},
        arrow/.style={-{Stealth}, thick, color=gray}
    ]
        % Layers
        \node[layer=primary] (client) {Couche Présentation (Jinja2 + Tailwind CSS + PWA)};
        \node[layer=secondary, below=0.3cm] (routes) {Couche Contrôleur (Flask Blueprints)};
        \node[layer=accent, below=0.3cm] (services) {Couche Service (Order, Payment, Analytics, QR, Loyalty)};
        \node[layer=orange, below=0.3cm] (models) {Couche Modèle (SQLAlchemy ORM)};
        \node[layer=teal, below=0.3cm] (db) {Couche Données (SQLite / PostgreSQL)};
        
        % WebSocket layer
        \node[layer=purple, right=1cm] at (routes.east) (websocket) {WebSocket (Flask-SocketIO)};
        
        % Arrows
        \draw[arrow] (client) -- (routes);
        \draw[arrow] (routes) -- (services);
        \draw[arrow] (services) -- (models);
        \draw[arrow] (models) -- (db);
        \draw[arrow, dashed] (websocket) -- (routes);
        \draw[arrow, dashed, bend left] (websocket) to (client);
    \end{tikzpicture}
    \caption{Architecture en couches de Tablii}
    \label{fig:architecture}
\end{figure}

\section{Architecture Multi-Tenant}

L'architecture multi-tenant de Tablii repose sur le principe d'\textbf{isolation logique} des données. Chaque restaurant (tenant) possède un identifiant unique (\texttt{restaurant\_id}) qui est utilisé comme clé étrangère sur toutes les tables tenant-spécifiques.

\begin{figure}[H]
    \centering
    \begin{tikzpicture}[
        box/.style={rectangle, rounded corners, draw=#1, fill=#1!5, minimum width=3cm, minimum height=1cm, font=\small},
        tenant/.style={rectangle, draw=#1, dashed, rounded corners, minimum width=5cm, minimum height=5cm}
    ]
        % Super Admin
        \node[box=primary, minimum width=4cm] (admin) at (0, 5) {\textbf{Super Admin}};
        
        % Tenant 1
        \node[tenant=secondary] (t1) at (-4, 0) {};
        \node[above=0.2cm] at (t1.north) {\textcolor{secondary}{Restaurant 1 (chez-ahmed)}};
        \node[box=secondary] (owner1) at (-4, 2) {Gérant};
        \node[box=accent] (cashier1) at (-6, -0.5) {Caissier};
        \node[box=orange] (kitchen1) at (-4, -0.5) {Cuisine};
        \node[box=purple] (waiter1) at (-2, -0.5) {Serveur};
        \node[box=teal] (data1) at (-4, -2.5) {\textit{Données isolées}};
        
        % Tenant 2
        \node[tenant=secondary] (t2) at (4, 0) {};
        \node[above=0.2cm] at (t2.north) {\textcolor{secondary}{Restaurant 2}};
        \node[box=secondary] (owner2) at (4, 2) {Gérant};
        \node[box=teal] (data2) at (4, -2.5) {\textit{Données isolées}};
        
        % Connections
        \draw[->, thick, primary] (admin) -- (t1);
        \draw[->, thick, primary] (admin) -- (t2);
        \draw[->, secondary] (owner1) -- (cashier1);
        \draw[->, secondary] (owner1) -- (kitchen1);
        \draw[->, secondary] (owner1) -- (waiter1);
        
        % Isolation note
        \node[draw=primary, fill=primary!5, rounded corners, text width=10cm, align=center] at (0, -5) {
            \textbf{Isolation :} Chaque restaurant a son propre \texttt{restaurant\_id}.\\
            Les requêtes sont filtrées automatiquement par \texttt{restaurant\_id}.
        };
    \end{tikzpicture}
    \caption{Architecture multi-tenant avec isolation des données}
    \label{fig:multitenant}
\end{figure}

\section{Modèle de Données}

La base de données de Tablii comprend \textbf{17 tables} organisées en domaines fonctionnels :

\begin{table}[H]
    \centering
    \begin{tabularx}{\textwidth}{l l l l}
        \toprule
        \rowcolor{primary!10}
        \textbf{Table} & \textbf{Clé Primaire} & \textbf{Clés Étrangères} & \textbf{Description} \\
        \midrule
        users & id & -- & Propriétaires \& super admins \\
        staff\_users & id & restaurant\_id $\to$ restaurants & Personnel de restaurant \\
        customers & id & -- & Clients finaux (par téléphone) \\
        restaurants & id & owner\_id $\to$ users & Restaurants (tenants) \\
        subscriptions & id & restaurant\_id $\to$ restaurants & Plans d'abonnement \\
        operating\_hours & id & restaurant\_id $\to$ restaurants & Horaires d'ouverture \\
        categories & id & restaurant\_id $\to$ restaurants & Catégories du menu \\
        menu\_items & id & category\_id $\to$ categories & Articles du menu \\
        customizations & id & menu\_item\_id $\to$ menu\_items & Groupes d'options \\
        custom\_options & id & customization\_id $\to$ customizations & Options individuelles \\
        tables\_ & id & restaurant\_id $\to$ restaurants & Tables physiques \\
        table\_sessions & id & table\_id $\to$ tables\_ & Sessions de dining \\
        orders & id & session\_id $\to$ table\_sessions & Commandes \\
        order\_items & id & order\_id $\to$ orders & Lignes de commande \\
        payment\_transactions & id & order\_id $\to$ orders & Transactions de paiement \\
        waiter\_calls & id & restaurant\_id $\to$ restaurants & Appels serveur \\
        reviews & id & order\_id $\to$ orders & Avis clients \\
        loyalty\_points & id & customer\_id $\to$ customers & Points de fidélité \\
        notifications & id & restaurant\_id $\to$ restaurants & Notifications staff \\
        \bottomrule
    \end{tabularx}
    \caption{Schéma complet de la base de données}
    \label{tab:schema}
\end{table}

\section{Flux de Statut des Commandes}

Le cycle de vie d'une commande suit une machine d'état stricte :

\begin{figure}[H]
    \centering
    \begin{tikzpicture}[
        state/.style={circle, draw=#1, fill=#1!10, minimum size=1.5cm, font=\small\bfseries, text=#1},
        trans/.style={-{Stealth}, thick, color=gray}
    ]
        \node[state=gray] (new) at (0, 0) {new};
        \node[state=secondary] (accepted) at (4, 0) {accepted};
        \node[state=accent] (preparing) at (8, 0) {preparing};
        \node[state=orange] (ready) at (12, 0) {ready};
        \node[state=purple] (served) at (12, -3) {served};
        \node[state=teal] (completed) at (8, -3) {completed};
        \node[state=primary] (rejected) at (4, -3) {rejected};
        
        \draw[trans] (new) -- node[above, font=\small] {Caissier accepte} (accepted);
        \draw[trans] (accepted) -- node[above, font=\small] {Cuisine commence} (preparing);
        \draw[trans] (preparing) -- node[above, font=\small] {Cuisine termine} (ready);
        \draw[trans] (ready) -- node[right, font=\small] {Serveur sert} (served);
        \draw[trans] (served) -- node[below, font=\small] {Paiement confirmé} (completed);
        \draw[trans] (new) -- node[left, font=\small] {Caissier refuse} (rejected);
        
        % Start
        \draw[trans] (-2, 0) -- node[above, font=\small] {Client commande} (new);
    \end{tikzpicture}
    \caption{Machine d'état des commandes}
    \label{fig:order-flow}
\end{figure}

\section{Diagramme de Séquence -- Passer une Commande}

\begin{figure}[H]
    \centering
    \begin{tikzpicture}[
        lifeline/.style={rectangle, draw, fill=lightgray, minimum width=2cm, minimum height=0.8cm, font=\small\bfseries},
        msg/.style={-{Stealth}, font=\small}
    ]
        % Lifelines
        \node[lifeline] (client) at (0, 6) {Client};
        \node[lifeline] (server) at (5, 6) {Flask Server};
        \node[lifeline] (db) at (10, 6) {Database};
        \node[lifeline] (ws) at (15, 6) {WebSocket};
        
        % Vertical lines
        \draw[dashed, gray] (client) -- ++(0, -12);
        \draw[dashed, gray] (server) -- ++(0, -12);
        \draw[dashed, gray] (db) -- ++(0, -12);
        \draw[dashed, gray] (ws) -- ++(0, -12);
        
        % Messages
        \draw[msg] (0, 5) -- node[above] {1. Scan QR code} (5, 4.8);
        \draw[msg] (5, 4.2) -- node[above] {2. GET /menu} (10, 4.0);
        \draw[msg, dashed] (10, 3.6) -- node[below] {3. Retourne menu} (5, 3.4);
        \draw[msg] (0, 2.8) -- node[above] {4. POST /order} (5, 2.6);
        \draw[msg] (5, 2.0) -- node[above] {5. INSERT order} (10, 1.8);
        \draw[msg] (5, 1.2) -- node[above] {6. emit('new\_order')} (15, 1.0);
        \draw[msg, dashed] (15, 0.6) -- node[below] {7. Notifie Caissier/Cuisine} (15, 0.4);
        \draw[msg, dashed] (10, 0) -- node[below] {8. Confirm} (5, -0.2);
        \draw[msg, dashed] (5, -0.8) -- node[below] {9. Confirmation client} (0, -1.0);
    \end{tikzpicture}
    \caption{Diagramme de séquence -- Processus de commande}
    \label{fig:sequence}
\end{figure}

\section{Architecture des Blueprints Flask}

L'application est structurée en \textbf{10 blueprints} Flask, chacun responsable d'un domaine fonctionnel :

\begin{figure}[H]
    \centering
    \begin{tikzpicture}[
        bp/.style={rectangle, rounded corners, draw=#1, fill=#1!5, minimum width=3cm, minimum height=0.8cm, font=\small},
        app/.style={rectangle, rounded corners, draw=primary, fill=primary!10, minimum width=3cm, minimum height=1cm, font=\bfseries}
    ]
        \node[app] (app) at (0, 5) {create\_app()};
        
        \node[bp=secondary] (auth) at (-8, 2) {auth};
        \node[bp=secondary] (landing) at (-4, 2) {landing};
        \node[bp=secondary] (onboarding) at (0, 2) {onboarding};
        \node[bp=secondary] (customer) at (4, 2) {customer};
        \node[bp=secondary] (dashboard) at (8, 2) {dashboard};
        
        \node[bp=accent] (cashier) at (-6, -1) {cashier};
        \node[bp=accent] (kitchen) at (-2, -1) {kitchen};
        \node[bp=accent] (waiter) at (2, -1) {waiter};
        \node[bp=accent] (admin) at (6, -1) {admin};
        \node[bp=accent] (api) at (10, -1) {api};
        
        \node[bp=orange, minimum width=4cm] (events) at (0, -4) {events (WebSocket)};
        
        \draw[->, thick, gray] (app) -- (auth);
        \draw[->, thick, gray] (app) -- (landing);
        \draw[->, thick, gray] (app) -- (onboarding);
        \draw[->, thick, gray] (app) -- (customer);
        \draw[->, thick, gray] (app) -- (dashboard);
        \draw[->, thick, gray] (app) -- (cashier);
        \draw[->, thick, gray] (app) -- (kitchen);
        \draw[->, thick, gray] (app) -- (waiter);
        \draw[->, thick, gray] (app) -- (admin);
        \draw[->, thick, gray] (app) -- (api);
        \draw[->, thick, gray, dashed] (app) -- (events);
    \end{tikzpicture}
    \caption{Architecture des Blueprints Flask}
    \label{fig:blueprints}
\end{figure}

\section{Services Métier}

\begin{table}[H]
    \centering
    \begin{tabularx}{\textwidth}{l l X}
        \toprule
        \rowcolor{accent!10}
        \textbf{Service} & \textbf{Fichier} & \textbf{Responsabilités} \\
        \midrule
        Order Service & order\_service.py & Création, mise à jour, calcul des totaux \\
        Payment Service & payment\_service.py & Intégration Flouci, webhooks, réconciliation \\
        Analytics Service & analytics\_service.py & Statistiques, heatmaps, tendances \\
        QR Service & qr\_service.py & Génération des QR codes pour les tables \\
        Loyalty Service & loyalty\_service.py & Points de fidélité, redemption \\
        Notification Service & notification\_service.py & Notifications in-app pour le staff \\
        Upload Service & upload\_service.py & Upload d'images (menu, logo, reviews) \\
        \bottomrule
    \end{tabularx}
    \caption{Services métier de l'application}
    \label{tab:services}
\end{table}

% ============================================================================
% CHAPITRE 4 : REALISATION ET IMPLÉMENTATION
% ============================================================================
\cleardoublepage
\chapter{Réalisation et Implémentation}

\section{Introduction}

Ce chapitre présente les aspects pratiques de l'implémentation de Tablii, incluant l'environnement de développement, la justification des choix technologiques, la structure du projet et les interfaces réalisées.

\section{Environnement Matériel et Logiciel}

\subsection{Environnement de Développement}

\begin{table}[H]
    \centering
    \begin{tabularx}{\textwidth}{l X}
        \toprule
        \rowcolor{primary!10}
        \textbf{Composant} & \textbf{Spécification} \\
        \midrule
        Système d'exploitation & Windows 11 / Linux Ubuntu 22.04 \\
        Processeur & Intel Core i5 / AMD Ryzen 5 ou supérieur \\
        Mémoire RAM & 8 Go minimum \\
        IDE & Visual Studio Code \\
        Contrôle de version & Git + GitHub \\
        \bottomrule
    \end{tabularx}
    \caption{Configuration matérielle et logicielle}
    \label{tab:env}
\end{table}

\subsection{Stack Technologique}

\begin{table}[H]
    \centering
    \begin{tabularx}{\textwidth}{l l X}
        \toprule
        \rowcolor{primary!10}
        \textbf{Couche} & \textbf{Technologie} & \textbf{Justification} \\
        \midrule
        Backend & Flask 3.0 + Python 3.11 & Lightweight, extensible, écosystème riche \\
        ORM & SQLAlchemy 2.0 & Type-safe, migrations Alembic \\
        Temps réel & Flask-SocketIO + eventlet & WebSocket natif, fallback polling \\
        Frontend & Jinja2 + Tailwind CSS 3.4 + Vanilla JS & Server-side rendering, utility-first CSS, pas de framework lourd \\
        Base de données & SQLite (dev) / PostgreSQL (prod) & Simplicité dev, robustesse prod \\
        Auth & Flask-Login + Werkzeug & Session-based, bcrypt hashing \\
        Paiement & Flouci API & Passerelle tunisienne \\
        Hébergement & Render.com & Deployment automatique, PostgreSQL managé \\
        \bottomrule
    \end{tabularx}
    \caption{Stack technologique de Tablii}
    \label{tab:techstack}
\end{table}

\subsection{Justification du Choix Vanilla JS + Tailwind CSS}

Le choix d'un frontend minimaliste en \textbf{Vanilla JavaScript} et \textbf{Tailwind CSS} est motivé par plusieurs facteurs critiques pour le contexte de la restauration en Tunisie :

\begin{itemize}[leftmargin=1.5cm, itemsep=0.3cm]
    \item \textbf{Performance extrême} : Sans la surcharge de React, Angular ou Vue.js, le temps de chargement initial est réduit de 60\%, ce qui est crucial pour les clients scannant un QR code sur un réseau mobile parfois instable.
    \item \textbf{Compatibilité maximale} : Vanilla JS fonctionne sur tous les navigateurs, y compris les versions anciennes encore utilisées sur les appareils d'entrée de gamme.
    \item \textbf{Tailwind CSS} : Le framework utility-first permet un design responsive et cohérent sans écrire de CSS personnalisé, tout en générant un fichier CSS minifié de moins de 15 Ko.
    \item \textbf{Maintenance simplifiée} : Pas de build step complexe pour le JavaScript, pas de gestion de dépendances npm côté frontend, ce qui réduit la complexité opérationnelle.
    \item \textbf{PWA native} : Le service worker et le manifest PWA sont implémentés directement sans dépendance externe, permettant l'installation sur l'écran d'accueil du smartphone.
\end{itemize}

\section{Structure du Projet}

\begin{figure}[H]
    \centering
    \begin{tikzpicture}[
        file/.style={font=\ttfamily\small, text=gray},
        dir/.style={font=\ttfamily\small\bfseries, text=secondary},
        root/.style={font=\ttfamily\bfseries, text=primary}
    ]
        \node[root] (root) at (0, 0) {Tablii/};
        \node[dir] (app) at (1, -0.8) {app/};
        \node[dir] (models) at (2.5, -1.6) {models/};
        \node[file] (m1) at (4, -2.2) {user.py, restaurant.py};
        \node[file] (m2) at (4, -2.7) {menu.py, order.py};
        \node[file] (m3) at (4, -3.2) {table.py, review.py};
        \node[dir] (routes) at (2.5, -4.2) {routes/};
        \node[file] (r1) at (4, -4.8) {auth.py, customer.py};
        \node[file] (r2) at (4, -5.3) {cashier.py, kitchen.py, waiter.py};
        \node[file] (r3) at (4, -5.8) {dashboard.py, admin.py, api.py};
        \node[dir] (services) at (2.5, -6.8) {services/};
        \node[file] (s1) at (4, -7.4) {order\_service.py, payment\_service.py};
        \node[file] (s2) at (4, -7.9) {analytics\_service.py, qr\_service.py};
        \node[dir] (events) at (2.5, -8.9) {events/};
        \node[file] (e1) at (4, -9.5) {order\_events.py, kitchen\_events.py};
        \node[dir] (templates) at (2.5, -10.5) {templates/};
        \node[dir] (static) at (2.5, -11.2) {static/};
        \node[file] (conf) at (1, -12) {config.py, run.py};
        \node[file] (req) at (1, -12.5) {requirements.txt};
        \node[file] (env) at (1, -13) {.env.example};
        \node[file] (seed) at (1, -13.5) {seed.py};
        \node[file] (tests) at (1, -14) {tests/};
        
        \draw[gray!30] (root) -- (app);
        \draw[gray!30] (app) -- (models);
        \draw[gray!30] (models) -- (m1);
        \draw[gray!30] (app) -- (routes);
        \draw[gray!30] (routes) -- (r1);
        \draw[gray!30] (app) -- (services);
        \draw[gray!30] (services) -- (s1);
        \draw[gray!30] (app) -- (events);
        \draw[gray!30] (events) -- (e1);
        \draw[gray!30] (app) -- (templates);
        \draw[gray!30] (app) -- (static);
        \draw[gray!30] (root) -- (conf);
    \end{tikzpicture}
    \caption{Structure arborescente du projet}
    \label{fig:structure}
\end{figure}

\section{Application Factory}

Le pattern Factory de Flask permet une initialisation modulaire de l'application :

\begin{lstlisting}[caption={Application Factory -- \texttt{app/\_\_init\_\_.py}}]
def create_app(config_name=None):
    """Application factory."""
    if config_name is None:
        config_name = os.environ.get('FLASK_ENV', 'development')

    app = Flask(__name__)
    app.config.from_object(config_by_name[config_name])

    # Initialize extensions
    db.init_app(app)
    migrate.init_app(app, db)
    login_manager.init_app(app)
    socketio.init_app(app, async_mode='threading')
    csrf.init_app(app)

    # Register blueprints
    register_blueprints(app)
    
    # Dual user loader (User + StaffUser)
    @login_manager.user_loader
    def load_user(user_id):
        if user_id.startswith('user_'):
            return db.session.get(User, int(user_id[5:]))
        elif user_id.startswith('staff_'):
            return db.session.get(StaffUser, int(user_id[6:]))
        return None

    return app
\end{lstlisting}

\section{Modèle de Commande}

\begin{lstlisting}[caption={Modèle Order -- \texttt{app/models/order.py}}]
class Order(db.Model):
    """A customer order placed at a table."""
    __tablename__ = 'orders'

    id = db.Column(db.Integer, primary_key=True)
    session_id = db.Column(db.Integer, db.ForeignKey('table_sessions.id'))
    restaurant_id = db.Column(db.Integer, db.ForeignKey('restaurants.id'), nullable=False)
    table_id = db.Column(db.Integer, db.ForeignKey('tables_.id'))
    customer_id = db.Column(db.Integer, db.ForeignKey('customers.id'))
    order_number = db.Column(db.String(10), nullable=False)
    status = db.Column(db.String(20), default='new')
    payment_method = db.Column(db.String(20))
    payment_status = db.Column(db.String(20), default='pending')
    subtotal = db.Column(db.Float, default=0.0)
    tax_amount = db.Column(db.Float, default=0.0)
    total_amount = db.Column(db.Float, default=0.0)
    is_gift = db.Column(db.Boolean, default=False)
    created_at = db.Column(db.DateTime, default=datetime.utcnow)
    accepted_at = db.Column(db.DateTime, nullable=True)
    preparing_at = db.Column(db.DateTime, nullable=True)
    ready_at = db.Column(db.DateTime, nullable=True)
    served_at = db.Column(db.DateTime, nullable=True)
    completed_at = db.Column(db.DateTime, nullable=True)
\end{lstlisting}

\section{WebSocket -- Temps Réel}

Les événements WebSocket permettent la communication en temps réel entre les différentes interfaces :

\begin{lstlisting}[caption={WebSocket Events -- \texttt{app/events/order\_events.py}}]
@socketio.on('new_order')
def handle_new_order(data):
    """Broadcast new order to cashier and kitchen."""
    restaurant_id = data.get('restaurant_id')
    emit('order_created', data, room=f'restaurant_{restaurant_id}')

@socketio.on('update_order_status')
def handle_status_update(data):
    """Broadcast order status change to all staff."""
    order_id = data.get('order_id')
    status = data.get('status')
    emit('status_changed', {
        'order_id': order_id,
        'status': status
    }, room=f'order_{order_id}')
\end{lstlisting}

\section{Interfaces Réalisées}

\subsection{Interface Client -- Menu QR Code}

L'interface client est accessible via le scan d'un QR code qui contient l'URL :
\texttt{https://tablii.com/menu/\{restaurant\_slug\}/table/\{table\_number\}}.

\begin{figure}[H]
    \centering
    \begin{tikzpicture}
        % Phone mockup
        \draw[rounded corners=20pt, thick, fill=dark] (-2, -4) rectangle (2, 4);
        \draw[rounded corners=16pt, fill=white] (-1.7, -3.5) rectangle (1.7, 3.5);
        
        % Header
        \fill[primary!10, rounded corners=16pt] (-1.7, 2.5) rectangle (1.7, 3.5);
        \node[font=\small\bfseries, text=primary] at (0, 3.2) {Chez Ahmed};
        \node[font=\tiny, text=gray] at (0, 2.8) {Table 5 -- Menu};
        
        % Menu items
        \node[font=\tiny\bfseries, anchor=west] at (-1.4, 2.2) {Entrées};
        \draw[rounded corners=4pt, fill=lightgray] (-1.5, 1.2) rectangle (1.5, 2.0);
        \node[font=\tiny, anchor=west] at (-1.3, 1.7) {Brik à l'oeuf};
        \node[font=\tiny, anchor=east] at (1.3, 1.7) {3.500 TND};
        
        \draw[rounded corners=4pt, fill=lightgray] (-1.5, 0.4) rectangle (1.5, 1.2);
        \node[font=\tiny, anchor=west] at (-1.3, 0.9) {Salade méchouia};
        \node[font=\tiny, anchor=east] at (1.3, 0.9) {4.000 TND};
        
        \node[font=\tiny\bfseries, anchor=west] at (-1.4, -0.2) {Plats Principaux};
        \draw[rounded corners=4pt, fill=primary!5] (-1.5, -1.0) rectangle (1.5, -0.2);
        \node[font=\tiny, anchor=west] at (-1.3, -0.5) {Couscroyal};
        \node[font=\tiny, anchor=east] at (1.3, -0.5) {12.000 TND};
        
        \draw[rounded corners=4pt, fill=lightgray] (-1.5, -1.8) rectangle (1.5, -1.0);
        \node[font=\tiny, anchor=west] at (-1.3, -1.3) {Tajine poulet};
        \node[font=\tiny, anchor=east] at (1.3, -1.3) {8.500 TND};
        
        % Cart button
        \fill[primary, rounded corners=8pt] (-1.2, -3.0) rectangle (1.2, -2.4);
        \node[white, font=\tiny\bfseries] at (0, -2.7) {Voir le Panier (2)};
        
        % Label
        \node[font=\small\bfseries, text=dark] at (0, -4.8) {Interface Client (Mobile)};
    \end{tikzpicture}
    \caption{Mockup de l'interface client mobile}
    \label{fig:client-mockup}
\end{figure}

\subsection{Interface Cuisine -- Kitchen Display System}

\begin{figure}[H]
    \centering
    \begin{tikzpicture}
        % Screen mockup
        \draw[rounded corners=8pt, thick, fill=dark] (-5, -3) rectangle (5, 3);
        \draw[rounded corners=4pt, fill=white] (-4.7, -2.7) rectangle (4.7, 2.7);
        
        % Header
        \fill[accent!10] (-4.7, 2.0) rectangle (4.7, 2.7);
        \node[font=\small\bfseries, text=accent, anchor=west] at (-4.5, 2.35) {Cuisine -- Commandes en cours};
        \node[font=\tiny, text=gray, anchor=east] at (4.5, 2.35) {\textbullet\ Connecté};
        
        % Order cards
        \draw[rounded corners=4pt, fill=secondary!10, draw=secondary] (-4.5, 0.8) rectangle (-0.5, 1.8);
        \node[font=\tiny\bfseries, text=secondary, anchor=west] at (-4.3, 1.5) {Commande \#1042};
        \node[font=\tiny, anchor=west] at (-4.3, 1.2) {2x Couscroyal, 1x Brik};
        \node[font=\tiny, anchor=west] at (-4.3, 0.9) {Table 5 -- il y a 2 min};
        
        \draw[rounded corners=4pt, fill=accent!10, draw=accent] (-0.2, 0.8) rectangle (3.8, 1.8);
        \node[font=\tiny\bfseries, text=accent, anchor=west] at (0.0, 1.5) {Commande \#1041};
        \node[font=\tiny, anchor=west] at (0.0, 1.2) {1x Tajine, 1x Salade};
        \node[font=\tiny, anchor=west] at (0.0, 0.9) {Table 3 -- En préparation};
        
        \draw[rounded corners=4pt, fill=orange!10, draw=orange] (-4.5, -1.0) rectangle (-0.5, 0.0);
        \node[font=\tiny\bfseries, text=orange, anchor=west] at (-4.3, -0.3) {Commande \#1040};
        \node[font=\tiny, anchor=west] at (-4.3, -0.6) {3x Grillade mixte};
        \node[font=\tiny, anchor=west] at (-4.3, -0.9) {Table 8 -- Prêt};
        
        % Label
        \node[font=\small\bfseries, text=dark] at (0, -3.5) {Interface Cuisine (Kitchen Display)};
    \end{tikzpicture}
    \caption{Mockup de l'interface cuisine}
    \label{fig:kitchen-mockup}
\end{figure}

\subsection{Dashboard Analytics}

Le tableau de bord gérant fournit des statistiques en temps réel :

\begin{figure}[H]
    \centering
    \begin{tikzpicture}
        % Dashboard mockup
        \draw[rounded corners=8pt, thick, fill=dark] (-5, -3.5) rectangle (5, 3.5);
        \draw[rounded corners=4pt, fill=white] (-4.7, -3.2) rectangle (4.7, 3.2);
        
        % Header
        \fill[primary!5] (-4.7, 2.5) rectangle (4.7, 3.2);
        \node[font=\small\bfseries, text=primary, anchor=west] at (-4.5, 2.85) {Dashboard -- Chez Ahmed};
        
        % KPI Cards
        \draw[rounded corners=4pt, fill=secondary!10] (-4.5, 1.5) rectangle (-1.5, 2.3);
        \node[font=\tiny, text=gray, anchor=west] at (-4.3, 2.1) {Revenu Aujourd'hui};
        \node[font=\small\bfseries, text=secondary, anchor=west] at (-4.3, 1.7) {1,250 TND};
        
        \draw[rounded corners=4pt, fill=accent!10] (-1.2, 1.5) rectangle (1.8, 2.3);
        \node[font=\tiny, text=gray, anchor=west] at (-1.0, 2.1) {Commandes};
        \node[font=\small\bfseries, text=accent, anchor=west] at (-1.0, 1.7) {47};
        
        \draw[rounded corners=4pt, fill=orange!10] (2.1, 1.5) rectangle (4.5, 2.3);
        \node[font=\tiny, text=gray, anchor=west] at (2.3, 2.1) {Ticket Moyen};
        \node[font=\small\bfseries, text=orange, anchor=west] at (2.3, 1.7) {26.6 TND};
        
        % Chart placeholder
        \draw[rounded corners=4pt, fill=lightgray] (-4.5, -1.0) rectangle (1.8, 1.3);
        \node[font=\tiny, text=gray] at (-1.35, 0.15) {Graphique Revenus (7 jours)};
        
        % Heatmap placeholder
        \draw[rounded corners=4pt, fill=lightgray] (2.1, -1.0) rectangle (4.5, 1.3);
        \node[font=\tiny, text=gray] at (3.3, 0.15) {Heatmap Heures};
        
        % Popular items
        \draw[rounded corners=4pt, fill=lightgray] (-4.5, -2.8) rectangle (4.5, -1.2);
        \node[font=\tiny, text=gray] at (0, -2.0) {Articles Populaires : 1. Couscroyal (23)  2. Brik (18)  3. Tajine (15)};
        
        % Label
        \node[font=\small\bfseries, text=dark] at (0, -4.0) {Interface Dashboard Gérant};
    \end{tikzpicture}
    \caption{Mockup du tableau de bord analytics}
    \label{fig:dashboard-mockup}
\end{figure}

\section{Mode Ramadan}

Tablii implémente un mode Ramadan spécial qui adapte le menu dynamiquement :

\begin{figure}[H]
    \centering
    \begin{tikzpicture}[
        box/.style={rectangle, rounded corners, draw=#1, fill=#1!5, minimum width=4cm, minimum height=1cm, font=\small}
    ]
        \node[box=primary] (mode) at (0, 3) {Mode Ramadan Activé};
        
        \node[box=accent] (iftar) at (-3, 0.5) {Iftar (18h30 -- 00h00)};
        \node[box=orange] (suhoor) at (3, 0.5) {Suhoor (00h00 -- 04h30)};
        
        \node[box=gray] (normal) at (0, -2) {Menu Normal (reste de la journée)};
        
        \draw[->, thick, primary] (mode) -- (iftar);
        \draw[->, thick, primary] (mode) -- (suhoor);
        \draw[->, thick, gray, dashed] (mode) -- (normal);
        
        \node[draw=primary, fill=primary!5, rounded corners, text width=8cm, align=center] at (0, -4) {
            \textbf{Logique :} Selon l'heure actuelle, le système affiche\\
            automatiquement les catégories Iftar ou Suhoor.
        };
    \end{tikzpicture}
    \caption{Fonctionnement du mode Ramadan}
    \label{fig:ramadan}
\end{figure}

\section{Intégration de Paiement Flouci}

\begin{lstlisting}[caption={Service de Paiement -- \texttt{app/services/payment\_service.py}}]
class PaymentService:
    """Flouci payment gateway integration."""
    
    def create_payment(self, order):
        """Create a payment session via Flouci API."""
        response = requests.post(
            'https://developers.flouci.com/api/generate_payment',
            headers={'app_token': FLOUCI_APP_TOKEN},
            json={
                'amount': order.total_amount,
                'currency': 'TND',
                'success_link': url_for('payment.success', order_id=order.id),
                'fail_link': url_for('payment.fail', order_id=order.id)
            }
        )
        return response.json()
    
    def verify_payment(self, payment_id):
        """Verify payment status via webhook or polling."""
        response = requests.get(
            f'https://developers.flouci.com/api/payments/{payment_id}',
            headers={'app_token': FLOUCI_APP_TOKEN}
        )
        return response.json()
\end{lstlisting}

\section{Déploiement sur Render.com}

\begin{figure}[H]
    \centering
    \begin{tikzpicture}[
        step/.style={rectangle, rounded corners, draw=#1, fill=#1!5, minimum width=3cm, minimum height=1cm, font=\small},
        arrow/.style={-{Stealth}, thick, color=gray}
    ]
        \node[step=secondary] (git) at (0, 4) {Push sur GitHub};
        \node[step=accent] (render) at (0, 2) {Render détecte render.yaml};
        \node[step=orange] (build) at (0, 0) {Build : pip install + Tailwind};
        \node[step=purple] (db) at (-3, -2) {PostgreSQL managé};
        \node[step=teal] (migrate) at (3, -2) {flask db upgrade};
        \node[step=primary] (deploy) at (0, -4) {Application en ligne};
        
        \draw[arrow] (git) -- (render);
        \draw[arrow] (render) -- (build);
        \draw[arrow] (build) -- (db);
        \draw[arrow] (build) -- (migrate);
        \draw[arrow] (db) -- (deploy);
        \draw[arrow] (migrate) -- (deploy);
    \end{tikzpicture}
    \caption{Pipeline de déploiement sur Render.com}
    \label{fig:deploy}
\end{figure}

\section{Résultats et Métriques}

\begin{table}[H]
    \centering
    \begin{tabularx}{\textwidth}{l X c}
        \toprule
        \rowcolor{accent!10}
        \textbf{Métrie} & \textbf{Description} & \textbf{Valeur} \\
        \midrule
        Modèles de données & Tables implémentées & 17 \\
        Blueprints Flask & Modules de routes & 10 \\
        Services métier & Classes de service & 7 \\
        Templates Jinja2 & Pages HTML & 40+ \\
        Endpoints API & Routes REST & 30+ \\
        Événements WebSocket & Handlers temps réel & 8 \\
        Couverture de tests & Code couvert par les tests & 75\%+ \\
        Temps de réponse API & Moyenne des requêtes & $<$ 100ms \\
        Latence WebSocket & Notification temps réel & $<$ 50ms \\
        \bottomrule
    \end{tabularx}
    \caption{Métriques du projet}
    \label{tab:metrics}
\end{table}

% ============================================================================
% CONCLUSION GENERALE
% ============================================================================
\cleardoublepage
\chapter*{Conclusion Générale et Perspectives}
\addcontentsline{toc}{chapter}{Conclusion Générale et Perspectives}

\section*{Bilan du Projet}

Tablii est une plateforme SaaS multi-tenant complète pour la gestion de restaurants, implémentant un ensemble de fonctionnalités modernes répondant aux spécificités du marché tunisien :

\begin{itemize}[leftmargin=1.5cm, itemsep=0.3cm]
    \item \textbf{Commandes via QR Code} : Les clients scannent et commandent sans application
    \item \textbf{Temps réel} : WebSocket pour la synchronisation instantanée entre tous les rôles
    \item \textbf{Multi-tenant} : Isolation complète des données entre restaurants
    \item \textbf{Multilingue} : Support français, arabe (RTL) et anglais
    \item \textbf{Paiement en ligne} : Intégration Flouci pour le marché tunisien
    \item \textbf{Analytics} : Dashboard avec statistiques, heatmaps et tendances
    \item \textbf{Mode Ramadan} : Adaptation automatique du menu Iftar/Suhoor
    \item \textbf{PWA} : Application web progressive installable
    \item \textbf{Fidélisation} : Système de points de loyauté
    \item \textbf{Performance} : Frontend Vanilla JS + Tailwind CSS ultra-léger
\end{itemize}

Le projet a été réalisé en suivant les bonnes pratiques de développement : pattern Factory, blueprints Flask, SQLAlchemy ORM, tests unitaires avec pytest, et déploiement automatisé sur Render.com.

\section*{Perspectives d'Évolution}

Plusieurs axes d'amélioration et d'évolution sont envisageables :

\begin{enumerate}[leftmargin=1.5cm, itemsep=0.3cm]
    \item \textbf{Application mobile native} : Développer une app iOS/Android avec React Native pour une expérience encore plus fluide
    \item \textbf{IA prédictive} : Prédiction des pics d'activité et recommandations de stock basées sur l'historique des commandes
    \item \textbf{Multi-restaurants} : Support des chaînes de restaurants avec vue consolidée pour les propriétaires de plusieurs établissements
    \item \textbf{Intégrations livraison} : Connexion avec les plateformes de livraison tunisiennes (Glovo, Yassir, etc.)
    \item \textbf{Réservation} : Module de réservation de tables en ligne avec gestion des créneaux
    \item \textbf{Inventaire} : Gestion des stocks et alertes de rupture de stock automatiques
    \item \textbf{Multi-langues additionnelles} : Turc, espagnol, italien pour cibler le marché international
    \item \textbf{API GraphQL} : Alternative à l'API REST pour des requêtes plus flexibles et optimisées
\end{enumerate}

\begin{center}
    \vspace{1cm}
    \begin{tikzpicture}
        \fill[primary!10, rounded corners] (-5, -1) rectangle (5, 1);
        \draw[primary, line width=2pt, rounded corners] (-5, -1) rectangle (5, 1);
        \node at (0, 0.3) {\textcolor{primary}{\Large\bfseries Tablii}};
        \node at (0, -0.3) {\textcolor{dark}{\large Digitaliser l'expérience de restauration}};
    \end{tikzpicture}
\end{center}

% ============================================================================
% BIBLIOGRAPHIE
% ============================================================================
\cleardoublepage
\begin{thebibliography}{99}
\addcontentsline{toc}{chapter}{Bibliographie}

\bibitem{flask}
Grinberg, M. (2018). \textit{Flask Web Development: Developing Web Applications with Python} (2nd ed.). O'Reilly Media.

\bibitem{sqlalchemy}
Bayer, R. (2023). \textit{SQLAlchemy 2.0 Documentation}. \url{https://docs.sqlalchemy.org}

\bibitem{websocket}
Loreto, S. \& Saint-Andre, P. (2011). \textit{The WebSocket Protocol}. RFC 6455. IETF.

\bibitem{tailwind}
Reinink, A. (2023). \textit{Tailwind CSS Documentation}. \url{https://tailwindcss.com}

\bibitem{saas}
Miller, D. (2020). \textit{SaaS Architecture Patterns}. \textit{ACM Computing Surveys}, 53(4), 1-35.

\bibitem{pwa}
Google Developers. (2023). \textit{Progressive Web Apps}. \url{https://web.dev/progressive-web-apps/}

\bibitem{flouci}
Flouci. (2024). \textit{Flouci Payment API Documentation}. \url{https://developers.flouci.com}

\bibitem{render}
Render. (2024). \textit{Render Deployment Documentation}. \url{https://render.com/docs}

\bibitem{multitenant}
Abbott, J. \& Fisher, M. (2015). \textit{The Art of Scalability} (2nd ed.). Addison-Wesley.

\bibitem{designpatterns}
Gamma, E., Helm, R., Johnson, R. \& Vlissides, J. (1994). \textit{Design Patterns: Elements of Reusable Object-Oriented Software}. Addison-Wesley.

\end{thebibliography}

% ============================================================================
% ANNEXES
% ============================================================================
\cleardoublepage
\appendix
\chapter{Guide d'Installation Locale}

\section{Prérequis}

\begin{itemize}
    \item Python 3.11+
    \item Node.js 18+
    \item Git
\end{itemize}

\section{Étapes d'Installation}

\begin{enumerate}
    \item Cloner le repository :
    \begin{lstlisting}[language=bash]
git clone https://github.com/0xyo/Tablii.git
cd Tablii
    \end{lstlisting}
    
    \item Créer et activer l'environnement virtuel :
    \begin{lstlisting}[language=bash]
python -m venv .venv
.\.venv\Scripts\Activate.ps1
    \end{lstlisting}
    
    \item Installer les dépendances Python :
    \begin{lstlisting}[language=bash]
pip install -r requirements.txt
    \end{lstlisting}
    
    \item Installer les dépendances Node :
    \begin{lstlisting}[language=bash]
npm install
    \end{lstlisting}
    
    \item Configurer l'environnement :
    \begin{lstlisting}[language=bash]
Copy-Item .env.example .env
    \end{lstlisting}
    
    \item Builder le CSS :
    \begin{lstlisting}[language=bash]
npm run build:css
    \end{lstlisting}
    
    \item Initialiser la base de données :
    \begin{lstlisting}[language=bash]
flask db upgrade
python seed.py
    \end{lstlisting}
    
    \item Lancer l'application :
    \begin{lstlisting}[language=bash]
python run.py
    \end{lstlisting}
\end{enumerate}

\section{Comptes de Démonstration}

\begin{table}[H]
    \centering
    \begin{tabularx}{\textwidth}{l l l l}
        \toprule
        \rowcolor{primary!10}
        \textbf{Rôle} & \textbf{Email/Username} & \textbf{Mot de passe} & \textbf{Slug} \\
        \midrule
        Super Admin & superadmin@tablii.com & admin1234 & -- \\
        Gérant & owner@tablii.com & owner1234 & chez-ahmed \\
        Caissier & caisse1 & staff1234 & chez-ahmed \\
        Serveur & serveur1 & staff1234 & chez-ahmed \\
        Cuisine & cuisine1 & staff1234 & chez-ahmed \\
        \bottomrule
    \end{tabularx}
    \caption{Comptes de démonstration}
    \label{tab:demo-accounts}
\end{table}

\chapter{Variables d'Environnement}

\begin{table}[H]
    \centering
    \begin{tabularx}{\textwidth}{l X l}
        \toprule
        \rowcolor{primary!10}
        \textbf{Variable} & \textbf{Description} & \textbf{Valeur par défaut} \\
        \midrule
        \texttt{SECRET\_KEY} & Clé de session Flask & \texttt{dev-fallback-change-me} \\
        \texttt{DATABASE\_URL} & URI de connexion BDD & \texttt{sqlite:///dev.db} \\
        \texttt{FLASK\_ENV} & Environnement & \texttt{development} \\
        \texttt{FLOUCI\_APP\_TOKEN} & Token Flouci & -- \\
        \texttt{FLOUCI\_APP\_SECRET} & Secret Flouci & -- \\
        \texttt{MAX\_CONTENT\_LENGTH} & Taille max upload & \texttt{5242880} (5 MB) \\
        \texttt{SOCKETIO\_CORS\_ORIGINS} & Origines WebSocket & \texttt{*} \\
        \bottomrule
    \end{tabularx}
    \caption{Variables d'environnement}
    \label{tab:env-vars}
\end{table}

\chapter{API REST -- Référence Rapide}

\begin{table}[H]
    \centering
    \begin{tabularx}{\textwidth}{l l X}
        \toprule
        \rowcolor{primary!10}
        \textbf{Méthode} & \textbf{Endpoint} & \textbf{Description} \\
        \midrule
        \texttt{GET} & \texttt{/api/v1/menu/\{slug\}} & Récupérer le menu d'un restaurant \\
        \texttt{POST} & \texttt{/api/v1/order} & Créer une nouvelle commande \\
        \texttt{GET} & \texttt{/api/v1/order/\{id\}} & Détails d'une commande \\
        \texttt{PUT} & \texttt{/api/v1/order/\{id\}/status} & Mettre à jour le statut \\
        \texttt{POST} & \texttt{/api/v1/payment} & Initier un paiement \\
        \texttt{GET} & \texttt{/api/v1/analytics} & Statistiques du restaurant \\
        \texttt{GET} & \texttt{/api/v1/reviews} & Avis du restaurant \\
        \texttt{POST} & \texttt{/api/v1/review} & Soumettre un avis \\
        \bottomrule
    \end{tabularx}
    \caption{Principaux endpoints API REST}
    \label{tab:api}
\end{table}

\end{document}